\chapter{Plano de Ensino (mesmo conteúdo do plano.tex)}

% Observação: este capítulo replica o conteúdo do arquivo aed/plano.tex
% como parte inicial do livro de notas de aula (professor).

\section{Visão Geral}
Esta disciplina tem dois objetivos complementares:
\begin{itemize}[leftmargin=*]
  \item \textbf{Unidade I (semanas 01--07):} construir a base de \textbf{AED} com \textbf{estatística prática} e \textbf{Python (Pandas)} para limpar, resumir e interpretar dados;
  \item \textbf{Unidade II (semanas 08--12):} transformar insights em \textbf{comunicação persuasiva} com \textbf{Tableau} e \textbf{storytelling}.
\end{itemize}

\begin{NoteBox}
\textbf{Regra do ritmo (aula noturna 19h--22h):} em semanas com trabalho, as \textbf{apresentações ocupam no máximo 1 hora}. O restante da aula continua com teoria + prática guiada.
\end{NoteBox}

\section{Calendário Letivo (12 semanas de aula)}
O calendário abaixo está alinhado ao arquivo \texttt{aed/calendario.md}. Semanas marcadas como \textit{não conta} (feriados/provas) não entram na contagem das 12 semanas.

\begin{center}
\renewcommand{\arraystretch}{1.15}
\begin{tabular}{@{}c p{3.2cm} p{4.2cm} p{5.7cm}@{}}
\toprule
\textbf{Semana} & \textbf{Período (Seg--Qua)} & \textbf{Tipo da semana} & \textbf{Observação} \\
\midrule
01 & 02/02 a 04/02 & Teoria + lab & Início das aulas \\
02 & 09/02 a 11/02 & Teoria + \textbf{Trabalho 1} & Entrega quinzenal \\
\midrule
\multicolumn{4}{@{}l}{\textit{14 a 18/02 --- Carnaval [F] (não conta semana)}}\\
\midrule
03 & 23/02 a 25/02 & Teoria + lab & Semana normal \\
04 & 02/03 a 04/03 & Teoria + \textbf{Trabalho 2} & Entrega quinzenal \\
05 & 09/03 a 11/03 & Teoria + lab & Semana normal \\
06 & 16/03 a 18/03 & Teoria + \textbf{Trabalho 3} & Entrega quinzenal \\
07 & 23/03 a 25/03 & Teoria + lab & Semana normal \\
\midrule
\multicolumn{4}{@{}l}{\textit{02 a 04/04 --- Semana Santa [F] (não conta semana)}}\\
\multicolumn{4}{@{}l}{\textit{06 a 14/04 --- Prova I Unidade (não conta semana)}}\\
\midrule
08 & 20/04 & Teoria + \textbf{Trabalho 4} & Aula apenas na segunda (21/04 é feriado) \\
09 & 27/04 a 29/04 & Teoria + lab & Semana normal \\
10 & 04/05 a 06/05 & Teoria + \textbf{Trabalho 5} & 06/05 é feriado (Vitória) \\
11 & 11/05 a 13/05 & Teoria + lab & Semana normal \\
12 & 18/05 a 20/05 & Teoria + \textbf{Trabalho 6} & Entrega quinzenal (encerramento do conteúdo) \\
\midrule
\multicolumn{4}{@{}l}{\textit{25/05 a 01/06 --- Prova II Unidade (não conta semana)}}\\
\bottomrule
\end{tabular}
\end{center}

\section{Avaliação (visão operacional)}
\begin{itemize}[leftmargin=*]
  \item \textbf{Trabalhos quinzenais (T1--T6):} entregas práticas curtas, com apresentação e checklist.
  \item \textbf{Prova I Unidade:} janela 06/04 a 14/04 (conforme calendário institucional).
  \item \textbf{Prova II Unidade:} janela 25/05 a 01/06.
\end{itemize}

\begin{NoteBox}
\textbf{Formato padrão de entrega (em todos os trabalhos):}
\begin{itemize}[leftmargin=*]
  \item \textbf{Python}: 1 notebook com ETL + checagens (tipos, ausências, duplicadas, filtros).
  \item \textbf{Dados}: 1 CSV ``limpo'' pronto para visualização.
  \item \textbf{Tableau}: 1 workbook (ou export) com visuais/dash.
  \item \textbf{Comunicação}: 1 slide (ou texto curto) com \textbf{Big Idea} + 3 bullets de evidências.
\end{itemize}
\end{NoteBox}

\section{Plano Semana a Semana (conteúdo + aprendizagem + entregas)}

\subsection*{Semana 01 (02/02--04/02) --- Do Caos à Clareza: AED e a \textit{Big Idea}}
\textbf{O que vai ter na semana:}
\begin{itemize}[leftmargin=*]
  \item Abertura com \textbf{gancho} e \textbf{Big Idea} do semestre: ``dados brutos são ruído; valor é extrair uma história que force uma decisão''.
  \item Diferença entre \textbf{análise exploratória} (cozinha) e \textbf{análise explanativa} (jantar).
  \item Primeiro \textbf{ETL com Python/Pandas}: leitura, inspeção, tipos, ausências, conversão de datas, remoção de colunas e exportação de CSV limpo.
  \item Primeiro contato com \textbf{Tableau}: conectar CSV e gerar visuais rápidos.
  \item Design básico: \textbf{Teste do Relance (3 segundos)} e remoção de ``lixo visual''.
\end{itemize}

\textbf{O aluno aprende a:}
\begin{itemize}[leftmargin=*]
  \item Diferenciar \textbf{descobrir} insight (AED) de \textbf{comunicar} insight (storytelling).
  \item Rodar um ETL simples e produzir um dataset ``pronto para viz''.
  \item Construir o primeiro visual no Tableau e avaliar clareza pelo Teste do Relance.
\end{itemize}

\begin{SolvedBox}
\textbf{Preparação para o Trabalho 1 (entrega na Semana 02):}
\begin{itemize}[leftmargin=*]
  \item Escolher um dataset simples (Kaggle ou base sugerida pelo professor).
  \item Definir 1 pergunta de negócio e 1 Big Idea.
\end{itemize}
\end{SolvedBox}

\subsection*{Semana 02 (09/02--11/02) --- Trabalho 1 (apresentações) + Fundamentos de dados retangulares}
\textbf{Formato da aula:}
\begin{itemize}[leftmargin=*]
  \item \textbf{Até 1 hora}: apresentações relâmpago do T1 (2--3 min por aluno/grupo).
  \item \textbf{Restante}: teoria + prática guiada.
\end{itemize}

\textbf{Teoria/Prática da semana:}
\begin{itemize}[leftmargin=*]
  \item \textbf{Dados retangulares}: linhas x colunas, variável x registro.
  \item \textbf{Dicionário de dados} e perguntas: ``o que cada coluna significa?''.
  \item Qualidade básica: ausências, duplicadas, categorias sujas, datas e textos.
\end{itemize}

\textbf{O aluno aprende a:}
\begin{itemize}[leftmargin=*]
  \item Formular uma pergunta e traduzir em colunas/transformações.
  \item Entregar uma história curta com evidência (Big Idea + visuais).
\end{itemize}

\begin{SolvedBox}
\textbf{Trabalho 1 (entrega na Semana 02):} ETL básico + \textbf{3 visuais} no Tableau + Big Idea (1 slide).
\end{SolvedBox}

\subsection*{Semana 03 (23/02--25/02) --- Estatística descritiva I: medidas de posição}
\textbf{O que vai ter na semana:}
\begin{itemize}[leftmargin=*]
  \item \textbf{Média, mediana, moda} e quando usar cada uma.
  \item \textbf{Resumo com Pandas}: \texttt{describe}, \texttt{value\_counts}, \texttt{groupby}.
  \item Leitura crítica: por que ``média'' pode mentir (assimetria/outliers).
\end{itemize}

\textbf{O aluno aprende a:}
\begin{itemize}[leftmargin=*]
  \item Resumir uma coluna numérica e uma coluna categórica.
  \item Explicar em linguagem de negócio o que um resumo indica.
\end{itemize}

\subsection*{Semana 04 (02/03--04/03) --- Trabalho 2 (apresentações) + Estatística descritiva II: variabilidade}
\textbf{Formato da aula:}
\begin{itemize}[leftmargin=*]
  \item \textbf{Até 1 hora}: apresentações do T2.
  \item \textbf{Restante}: teoria + prática (boxplot, IQR, dispersão).
\end{itemize}

\textbf{Teoria/Prática da semana:}
\begin{itemize}[leftmargin=*]
  \item \textbf{Variância, desvio-padrão, amplitude, IQR} e leitura prática.
  \item \textbf{Outliers}: por que aparecem e como investigar antes de remover.
\end{itemize}

\textbf{O aluno aprende a:}
\begin{itemize}[leftmargin=*]
  \item Comparar dois grupos (ex.: por categoria) usando medidas de posição e dispersão.
  \item Justificar escolhas (mediana vs média; IQR vs DP) com base no dado.
\end{itemize}

\begin{SolvedBox}
\textbf{Trabalho 2 (entrega na Semana 04):} mini-relatório AED (posição + dispersão) + 2 visuais (boxplot + histograma) + Big Idea.
\end{SolvedBox}

\subsection*{Semana 05 (09/03--11/03) --- Distribuições e dados ausentes: o ``formato'' dos dados}
\textbf{O que vai ter na semana:}
\begin{itemize}[leftmargin=*]
  \item \textbf{Histogramas} e leitura de forma (simetria/assimetria; caudas).
  \item \textbf{Valores ausentes}: tipos de missing e estratégias (imputar, remover, sinalizar).
  \item \textbf{Regras práticas} para preparar dados para o Tableau (tipos, datas, categorias).
\end{itemize}

\textbf{O aluno aprende a:}
\begin{itemize}[leftmargin=*]
  \item Identificar distribuições problemáticas e propor transformações simples.
  \item Tratar missing sem ``matar'' o significado do dado.
\end{itemize}

\subsection*{Semana 06 (16/03--18/03) --- Trabalho 3 (apresentações) + Correlação introdutória}
\textbf{Formato da aula:}
\begin{itemize}[leftmargin=*]
  \item \textbf{Até 1 hora}: apresentações do T3.
  \item \textbf{Restante}: teoria + prática (correlação e gráficos bivariados).
\end{itemize}

\textbf{Teoria/Prática da semana:}
\begin{itemize}[leftmargin=*]
  \item \textbf{Dispersão (scatter)} e o que \textbf{correlação} mede (e o que não mede).
  \item \textbf{Segmentação} por categoria (cor por classe, facetas simples).
\end{itemize}

\textbf{O aluno aprende a:}
\begin{itemize}[leftmargin=*]
  \item Investigar relações entre duas variáveis sem cair em conclusões causais.
  \item Contar uma história com evidência bivariada (2 visuais coerentes).
\end{itemize}

\begin{SolvedBox}
\textbf{Trabalho 3 (entrega na Semana 06):} diagnóstico (missing + outliers + distribuição) + 1 relação bivariada (scatter) + Big Idea.
\end{SolvedBox}

\subsection*{Semana 07 (23/03--25/03) --- Variáveis categóricas e revisão para Prova I}
\textbf{O que vai ter na semana:}
\begin{itemize}[leftmargin=*]
  \item \textbf{Categóricas}: proporções, tabelas de frequência e gráficos de barras.
  \item \textbf{Boas práticas} de escolha de gráfico (evitar pizza quando não faz sentido).
  \item Revisão guiada dos conceitos da Unidade I (ETL, resumos, distribuições, correlação).
\end{itemize}

\textbf{O aluno aprende a:}
\begin{itemize}[leftmargin=*]
  \item Comparar categorias com clareza e consistência visual.
  \item Consolidar um ``kit'' de AED para a Prova I.
\end{itemize}

\subsection*{Semana 08 (20/04) --- Trabalho 4 (apresentações) + Storytelling: da análise à decisão}
\textbf{Observação:} aula concentrada (apenas segunda-feira).

\textbf{Formato da aula:}
\begin{itemize}[leftmargin=*]
  \item \textbf{Até 1 hora}: apresentações do T4.
  \item \textbf{Restante}: storytelling aplicado (estrutura e mensagem).
\end{itemize}

\textbf{Teoria/Prática da semana:}
\begin{itemize}[leftmargin=*]
  \item \textbf{AED vs explanativa} (retomada): o que entra e o que não entra numa apresentação.
  \item \textbf{Big Idea} e \textbf{pergunta de decisão}: ``o que eu quero que a audiência faça?''
  \item Primeiro \textbf{storyboard} (3 a 5 telas) para um dashboard.
\end{itemize}

\textbf{O aluno aprende a:}
\begin{itemize}[leftmargin=*]
  \item Transformar achados em uma linha narrativa curta.
  \item Planejar um dashboard antes de construir.
\end{itemize}

\begin{SolvedBox}
\textbf{Trabalho 4 (entrega na Semana 08):} roteiro curto (Big Idea + 3 evidências) + esboço de dashboard (wireframe) + 1 visual no Tableau.
\end{SolvedBox}

\subsection*{Semana 09 (27/04--29/04) --- Design de informação: clareza, hierarquia e consistência}
\textbf{O que vai ter na semana:}
\begin{itemize}[leftmargin=*]
  \item \textbf{Teste do Relance} em profundidade: título que diz a mensagem, não o tipo de gráfico.
  \item \textbf{Hierarquia visual}: cor, tamanho, espaço em branco, alinhamento.
  \item Padrões práticos para dashboards: filtros, contexto, explicações curtas.
\end{itemize}

\textbf{O aluno aprende a:}
\begin{itemize}[leftmargin=*]
  \item Melhorar um visual ``ruim'' com pequenas decisões de design.
  \item Escrever títulos orientados a insight (mensagem).
\end{itemize}

\subsection*{Semana 10 (04/05--06/05) --- Trabalho 5 (apresentações) + Tableau interativo}
\textbf{Observação:} 06/05 é feriado; ajustar a prática para caber em 2 encontros.

\textbf{Formato da aula:}
\begin{itemize}[leftmargin=*]
  \item \textbf{Até 1 hora}: apresentações do T5.
  \item \textbf{Restante}: Tableau avançado para interatividade.
\end{itemize}

\textbf{Teoria/Prática da semana:}
\begin{itemize}[leftmargin=*]
  \item Filtros, ações, parâmetros e interação guiada.
  \item Público-alvo: o que o usuário precisa explorar sozinho.
\end{itemize}

\textbf{O aluno aprende a:}
\begin{itemize}[leftmargin=*]
  \item Criar pelo menos 1 mecanismo de exploração (filtro/parâmetro/ação).
  \item Preservar clareza mesmo com interatividade.
\end{itemize}

\begin{SolvedBox}
\textbf{Trabalho 5 (entrega na Semana 10):} dashboard com 1 a 2 mecanismos de interatividade + título-mensagem + checklist do Teste do Relance.
\end{SolvedBox}

\subsection*{Semana 11 (11/05--13/05) --- Dashboards que ``forçam'' decisão}
\textbf{O que vai ter na semana:}
\begin{itemize}[leftmargin=*]
  \item Estrutura de dashboard: contexto, diagnóstico, recomendação.
  \item Consistência de métricas e definições (evitar ``cada um calcula de um jeito'').
  \item Preparar ``pitch'' de 3 minutos (história + decisão + próximo passo).
\end{itemize}

\textbf{O aluno aprende a:}
\begin{itemize}[leftmargin=*]
  \item Organizar telas para guiar o olhar e reduzir ruído.
  \item Conectar insight a recomendação (chamada para ação).
\end{itemize}

\subsection*{Semana 12 (18/05--20/05) --- Trabalho 6 (apresentações) + Fechamento do conteúdo}
\textbf{Formato da aula:}
\begin{itemize}[leftmargin=*]
  \item \textbf{Até 1 hora}: apresentações do T6 (pitch + demo do dashboard).
  \item \textbf{Restante}: refinamento, revisão e checklist para Prova II.
\end{itemize}

\textbf{O aluno aprende a:}
\begin{itemize}[leftmargin=*]
  \item Defender um dashboard como produto (para usuário e decisão).
  \item Fechar lacunas técnicas (dados) e de comunicação (mensagem).
\end{itemize}

\begin{SolvedBox}
\textbf{Trabalho 6 (entrega na Semana 12):} dashboard final + roteiro de apresentação (3 min) + checklist completo (dados limpos + Big Idea + interatividade).
\end{SolvedBox}

\section{Resumo das Entregas Quinzenais}
\begin{center}
\renewcommand{\arraystretch}{1.15}
\begin{tabular}{@{}c c p{10.5cm}@{}}
\toprule
\textbf{Entrega} & \textbf{Semana} & \textbf{Tema (artefatos principais)} \\
\midrule
T1 & 02 & ETL básico + 3 visuais + Big Idea \\
T2 & 04 & Posição + dispersão + boxplot/hist + Big Idea \\
T3 & 06 & Missing/outliers/distribuição + 1 relação bivariada \\
T4 & 08 & Storytelling + wireframe + 1 visual no Tableau \\
T5 & 10 & Dashboard interativo (filtro/parâmetro/ação) + título-mensagem \\
T6 & 12 & Dashboard final + pitch (3 min) + checklist \\
\bottomrule
\end{tabular}
\end{center}

% ===========================================================================
\section{Criando uma Biblioteca Python: da Ideia ao PyPI}
% ===========================================================================

Esta seção explica, passo a passo, como um conjunto de funções Python
evolui de um módulo local para um \textbf{pacote instalável mundialmente}
via \texttt{pip install}. É exatamente o que faremos com a biblioteca
\texttt{aedtools} ao longo do semestre.

\subsection*{Por que criar uma biblioteca?}

Sem biblioteca, cada notebook repete as mesmas funções. Com uma
biblioteca bem estruturada:
\begin{itemize}[leftmargin=*]
  \item \textbf{Reutilização:} qualquer notebook importa \texttt{from aedtools import calcular\_media}.
  \item \textbf{Versão única da verdade:} corrigiu o bug uma vez, todos os notebooks se beneficiam.
  \item \textbf{Documentação viva:} docstrings, tipos e exemplos ficam junto ao código.
  \item \textbf{Colaboração:} equipes inteiras usam o mesmo toolkit sem copiar arquivos.
\end{itemize}

\subsection*{Estrutura de pastas obrigatória}

O padrão moderno (recomendado pelo Python Packaging Authority) é o layout
\textbf{src}:

\begin{lstlisting}[language={}]
aedtools/                    <- raiz do projeto (repositório Git)
+-- src/
|   +-- aedtools/            <- pacote Python (nome real da lib)
|       +-- __init__.py      <- exporta a API pública
|       +-- estatistica.py   <- medidas de posição e dispersão
|       +-- limpeza.py       <- nulos e imputação
|       +-- transformacao.py <- normalização, encoding, ranking
|       +-- correlacao.py    <- Pearson, métricas de erro, IC
|       +-- visualizacao.py  <- gráficos (matplotlib / seaborn)
+-- tests/
|   +-- __init__.py
|   +-- test_estatistica.py  <- testes unitários (pytest)
|   +-- test_limpeza.py
+-- pyproject.toml           <- metadados + dependências
+-- README.md                <- documentação mínima
+-- LICENSE                  <- MIT é a mais comum para código aberto
\end{lstlisting}

\begin{NoteBox}
Por que \texttt{src/}? Evita que o Python importe acidentalmente o código
do diretório de trabalho em vez do pacote instalado — um bug silencioso
muito comum em projetos sem essa separação.
\end{NoteBox}

\subsection*{O arquivo \texttt{pyproject.toml}}

Substitui o antigo \texttt{setup.py} e o \texttt{setup.cfg}. É o único
arquivo de configuração necessário para projetos modernos:

\begin{lstlisting}[language={}]
[build-system]
requires      = ["hatchling"]       # backend de build
build-backend = "hatchling.build"

[project]
name            = "aedtools"        # nome no PyPI (deve ser único)
version         = "0.1.0"           # semântica: MAJOR.MINOR.PATCH
description     = "Biblioteca de AED para a turma UNIFACOL SI-8"
readme          = "README.md"
license         = {file = "LICENSE"}
requires-python = ">=3.10"          # versão mínima do Python
authors         = [{name = "Turma AED 2026-1", email = "aed@unifacol.edu.br"}]
keywords        = ["AED", "estatística", "pandas", "visualização"]

# Dependências que serão instaladas junto com a lib
dependencies = [
  "pandas>=2.0",
  "numpy>=1.26",
  "matplotlib>=3.7",
  "seaborn>=0.13",
  "scipy>=1.11",
]

[project.urls]
Repository = "https://github.com/unifacol/aedtools"

[tool.hatch.build.targets.wheel]
packages = ["src/aedtools"]         # de onde o wheel é construído
\end{lstlisting}

\subsection*{O arquivo \texttt{\_\_init\_\_.py} — a API pública}

O \texttt{\_\_init\_\_.py} decide o que o usuário enxerga ao digitar
\texttt{from aedtools import ...}:

\begin{lstlisting}[language=Python]
# src/aedtools/__init__.py
# Importações que tornam "from aedtools import X" possível
# sem precisar saber em qual módulo X está definido.

from .estatistica import (
    calcular_media, calcular_mediana, calcular_moda,
    calcular_desvio_padrao, calcular_variancia, calcular_amplitude,
    calcular_quartis, calcular_iqr, detectar_outliers_iqr,
    calcular_skewness, calcular_kurtosis,
)
from .limpeza import (
    contar_nulos, proporcao_nulos,
    preencher_nulos_media, preencher_nulos_mediana,
    remover_linhas_nulas,
)
from .transformacao import (
    normalizar_min_max, padronizar_zscore,
    contar_valores_unicos, calcular_frequencia,
    converter_para_categorico, one_hot_encoding,
    soma_acumulada, media_movel,
    variacao_percentual, calcular_ranking,
    identificar_tipos_colunas,
)
from .correlacao import (
    calcular_correlacao, gerar_matriz_correlacao,
    calcular_r2, calcular_mae, calcular_mse, calcular_rmse,
    intervalo_confianca_media,
)
from .visualizacao import (
    plotar_histograma, plotar_boxplot, plotar_scatterplot,
    plotar_heatmap_correlacao,
    plotar_barras_categorias, plotar_pizza_categorias,
)

# Versão acessível por aedtools.__version__
__version__ = "0.1.0"
__all__ = [
    "calcular_media", "calcular_mediana", "calcular_moda",
    # ... (lista completa de nomes exportados)
]
\end{lstlisting}

\subsection*{Comandos de desenvolvimento — passo a passo}

\begin{lstlisting}[language={}]
# 1. Criar e ativar ambiente virtual isolado
python -m venv .venv
source .venv/bin/activate          # Linux/macOS
# .venv\Scripts\activate           # Windows

# 2. Instalar as ferramentas de build e testes
pip install hatchling build twine pytest

# 3. Instalar a própria lib em modo editável (editable install)
#    Qualquer mudança no código src/ reflete imediatamente
#    sem precisar reinstalar.
pip install -e .

# 4. Rodar os testes
pytest tests/ -v

# 5. Construir os artefatos distribuíveis
python -m build
#  Cria dois arquivos em dist/:
#    aedtools-0.1.0.tar.gz    <- source distribution (sdist)
#    aedtools-0.1.0-py3-none-any.whl <- wheel (instalação rápida)

# 6. Conferir o pacote antes de publicar
twine check dist/*

# 7. Publicar no TestPyPI (ambiente de testes — gratuito)
twine upload --repository testpypi dist/*
#  Testar instalação:
pip install --index-url https://test.pypi.org/simple/ aedtools

# 8. Publicar no PyPI real (qualquer pessoa no mundo instala)
twine upload dist/*
#  Qualquer pessoa no planeta pode então usar:
pip install aedtools
\end{lstlisting}

\begin{AttentionBox}
Para publicar no PyPI você precisa de uma conta em \texttt{pypi.org} e
configurar um \textbf{API token} (nunca username/senha em scripts).
O token é gerado em Configurações de Conta $\to$ API tokens.
Guarde-o em \texttt{\textasciitilde/.pypirc} ou como variável de ambiente
\texttt{TWINE\_PASSWORD}.
\end{AttentionBox}

\subsection*{Diferença entre \texttt{.tar.gz} (sdist) e \texttt{.whl} (wheel)}

\begin{center}
\renewcommand{\arraystretch}{1.2}
\begin{tabular}{@{}l p{5cm} p{5.5cm}@{}}
\toprule
& \textbf{.tar.gz (sdist)} & \textbf{.whl (wheel)} \\
\midrule
Conteúdo & Código-fonte completo & Código já pronto para instalação \\
Build no cliente & Sim (pode precisar compilar) & Não (apenas descompactar) \\
Instalação & Mais lenta & Muito mais rápida \\
Portabilidade & Total (qualquer OS) & Depende da plataforma se há C/Cython \\
Quando usar & Libs com extensões C & Libs puro-Python (como aedtools) \\
\bottomrule
\end{tabular}
\end{center}

\subsection*{Boas práticas de desenvolvimento de biblioteca}

\begin{enumerate}[leftmargin=*]
  \item \textbf{Docstrings completas:} use o padrão NumPy ou Google.
  Parâmetros, retorno, exceções e pelo menos 1 exemplo em \texttt{Examples}.
  \item \textbf{Type hints obrigatórios:} \texttt{def calcular\_media(df: pd.DataFrame, coluna: str) -> float}
  — permite IDE autocompletar para o usuário.
  \item \textbf{Nunca modificar o DataFrame original:} sempre faça
  \texttt{df = df.copy()} antes de qualquer mutação. Funções puras são previsíveis.
  \item \textbf{Validar entradas:} cheque se a coluna existe, se é numérica etc.
  Lance \texttt{ValueError} com mensagem clara — não deixe o erro do Pandas confundir o usuário.
  \item \textbf{Testes unitários para cada função:} use \texttt{pytest}.
  Cada função deve ter ao menos: caso normal, coluna com NaN, coluna vazia.
  \item \textbf{Versionamento semântico:} MAJOR.MINOR.PATCH.
  Bug fix $\to$ patch; nova função $\to$ minor; mudança de API quebra compatibilidade $\to$ major.
  \item \textbf{CHANGELOG.md:} registre o que mudou em cada versão.
  Usuários precisam saber o que a atualização traz.
  \item \textbf{CI/CD (GitHub Actions):} automatize o \texttt{pytest} a cada push.
  Não suba código que não passou nos testes.
\end{enumerate}

\subsection*{Como usar a biblioteca após instalação}

\begin{lstlisting}[language=Python]
# Instalação
# pip install aedtools

import pandas as pd
from aedtools import (
    calcular_media, calcular_mediana, detectar_outliers_iqr,
    plotar_boxplot, calcular_frequencia,
)

# Carregar dados
df = pd.read_csv("vendas.csv")

# Usar as funções da biblioteca
media   = calcular_media(df, "faturamento")
mediana = calcular_mediana(df, "faturamento")
outliers = detectar_outliers_iqr(df, "faturamento")

print(f"Média: R$ {media:,.2f}")
print(f"Mediana: R$ {mediana:,.2f}")
print(f"Outliers encontrados: {len(outliers)} registros")

freq = calcular_frequencia(df, "canal")
print(freq)

plotar_boxplot(df, "faturamento")
\end{lstlisting}

% ===========================================================================
\section{\texttt{aedtools} — A Biblioteca da Turma: 40 Funções Implementadas}
% ===========================================================================

O nome \textbf{aedtools} (\textit{Análise Exploratória de Dados — Tools})
foi escolhido por ser curto, sem acento, descritivo e disponível no PyPI.
Cada uma das 40 funções abaixo está implementada com:
comentários linha a linha, docstring completa, validação de entrada
e exemplo de uso.

% ---------------------------------------------------------------------------
\subsection{Módulo \texttt{estatistica.py} — Medidas de Posição, Dispersão e Forma}
% ---------------------------------------------------------------------------

\begin{lstlisting}[language=Python]
# src/aedtools/estatistica.py
"""
Medidas de posição (média, mediana, moda),
dispersão (variância, DP, amplitude, quartis, IQR)
e forma (skewness, kurtosis).
"""

import pandas as pd
import numpy as np
from scipy import stats


# ── Função 1 ────────────────────────────────────────────────────────────────
def calcular_media(df: pd.DataFrame, coluna: str) -> float:
    """
    Calcula a média aritmética de uma coluna numérica.

    Parameters
    ----------
    df     : pd.DataFrame  — dataset de entrada
    coluna : str           — nome da coluna numérica

    Returns
    -------
    float — média aritmética ignorando NaN

    Examples
    --------
    >>> calcular_media(df, "faturamento")
    5432.1
    """
    # Verifica se a coluna existe no DataFrame
    if coluna not in df.columns:
        raise ValueError(f"Coluna '{coluna}' não encontrada no DataFrame.")
    # Verifica se a coluna é numérica
    if not pd.api.types.is_numeric_dtype(df[coluna]):
        raise TypeError(f"Coluna '{coluna}' não é numérica.")
    # Calcula a média ignorando valores nulos (skipna=True é padrão)
    return float(df[coluna].mean())


# ── Função 2 ────────────────────────────────────────────────────────────────
def calcular_mediana(df: pd.DataFrame, coluna: str) -> float:
    """
    Calcula a mediana (percentil 50) de uma coluna numérica.
    Robusta a outliers: não é afetada por valores extremos.

    Parameters
    ----------
    df     : pd.DataFrame
    coluna : str

    Returns
    -------
    float — mediana da coluna
    """
    # Validações de entrada
    if coluna not in df.columns:
        raise ValueError(f"Coluna '{coluna}' não encontrada.")
    if not pd.api.types.is_numeric_dtype(df[coluna]):
        raise TypeError(f"Coluna '{coluna}' não é numérica.")
    # median() ignora NaN automaticamente
    return float(df[coluna].median())


# ── Função 3 ────────────────────────────────────────────────────────────────
def calcular_moda(df: pd.DataFrame, coluna: str) -> list:
    """
    Calcula a(s) moda(s) de uma coluna — funciona para numérica
    e categórica. Retorna lista pois pode haver múltiplas modas.

    Parameters
    ----------
    df     : pd.DataFrame
    coluna : str

    Returns
    -------
    list — lista com os valores mais frequentes
    """
    if coluna not in df.columns:
        raise ValueError(f"Coluna '{coluna}' não encontrada.")
    # mode() retorna Series com todos os valores de máxima frequência
    # tolist() converte para lista Python padrão
    return df[coluna].mode().tolist()


# ── Função 4 ────────────────────────────────────────────────────────────────
def calcular_desvio_padrao(
    df: pd.DataFrame, coluna: str, populacional: bool = False
) -> float:
    """
    Calcula o desvio padrão amostral (ddof=1) ou populacional (ddof=0).

    Parameters
    ----------
    df          : pd.DataFrame
    coluna      : str
    populacional: bool — se True usa N; se False usa N-1 (padrão amostral)

    Returns
    -------
    float — desvio padrão
    """
    if coluna not in df.columns:
        raise ValueError(f"Coluna '{coluna}' não encontrada.")
    if not pd.api.types.is_numeric_dtype(df[coluna]):
        raise TypeError(f"Coluna '{coluna}' não é numérica.")
    # ddof=0 -> divide por N (populacional)
    # ddof=1 -> divide por N-1 (amostral — correcão de Bessel)
    ddof = 0 if populacional else 1
    return float(df[coluna].std(ddof=ddof))


# ── Função 5 ────────────────────────────────────────────────────────────────
def calcular_variancia(
    df: pd.DataFrame, coluna: str, populacional: bool = False
) -> float:
    """
    Calcula a variância amostral (ddof=1) ou populacional (ddof=0).
    A variância é o quadrado do desvio padrão — sua unidade é o
    quadrado da unidade original, portanto menos interpretável diretamente.

    Parameters
    ----------
    df          : pd.DataFrame
    coluna      : str
    populacional: bool

    Returns
    -------
    float — variância
    """
    if coluna not in df.columns:
        raise ValueError(f"Coluna '{coluna}' não encontrada.")
    if not pd.api.types.is_numeric_dtype(df[coluna]):
        raise TypeError(f"Coluna '{coluna}' não é numérica.")
    ddof = 0 if populacional else 1
    # var() segue a mesma lógica do std() com ddof
    return float(df[coluna].var(ddof=ddof))


# ── Função 6 ────────────────────────────────────────────────────────────────
def calcular_amplitude(df: pd.DataFrame, coluna: str) -> float:
    """
    Calcula a amplitude total: diferença entre o valor máximo e mínimo.
    Muito sensível a outliers — use apenas como primeira aproximação.

    Parameters
    ----------
    df     : pd.DataFrame
    coluna : str

    Returns
    -------
    float — amplitude (max - min)
    """
    if coluna not in df.columns:
        raise ValueError(f"Coluna '{coluna}' não encontrada.")
    if not pd.api.types.is_numeric_dtype(df[coluna]):
        raise TypeError(f"Coluna '{coluna}' não é numérica.")
    # Ignora NaN ao calcular max e min
    return float(df[coluna].max() - df[coluna].min())


# ── Função 11 ───────────────────────────────────────────────────────────────
def calcular_quartis(df: pd.DataFrame, coluna: str) -> dict:
    """
    Calcula Q1 (25%), Q2/mediana (50%), Q3 (75%) e IQR de uma coluna.

    Parameters
    ----------
    df     : pd.DataFrame
    coluna : str

    Returns
    -------
    dict — chaves: "Q1", "Q2", "Q3", "IQR"
    """
    if coluna not in df.columns:
        raise ValueError(f"Coluna '{coluna}' não encontrada.")
    if not pd.api.types.is_numeric_dtype(df[coluna]):
        raise TypeError(f"Coluna '{coluna}' não é numérica.")
    serie = df[coluna].dropna()   # remove NaN antes do cálculo
    q1 = float(serie.quantile(0.25))
    q2 = float(serie.quantile(0.50))
    q3 = float(serie.quantile(0.75))
    iqr = q3 - q1                 # Intervalo Interquartil
    return {"Q1": q1, "Q2": q2, "Q3": q3, "IQR": iqr}


# ── Função 12 ───────────────────────────────────────────────────────────────
def calcular_iqr(df: pd.DataFrame, coluna: str) -> float:
    """
    Retorna apenas o IQR (Q3 - Q1) — versão simplificada de calcular_quartis.
    Útil quando somente o IQR é necessário (ex.: detecção de outliers).

    Parameters
    ----------
    df     : pd.DataFrame
    coluna : str

    Returns
    -------
    float — IQR
    """
    # Reutiliza calcular_quartis para evitar duplicação de lógica
    quartis = calcular_quartis(df, coluna)
    return quartis["IQR"]


# ── Função 13 ───────────────────────────────────────────────────────────────
def detectar_outliers_iqr(df: pd.DataFrame, coluna: str) -> pd.DataFrame:
    """
    Detecta outliers pela Regra de Tukey (1,5 × IQR).
    Retorna subconjunto do DataFrame com apenas as linhas outlier.

    Limites:
        Inferior = Q1 - 1.5 * IQR
        Superior = Q3 + 1.5 * IQR

    Parameters
    ----------
    df     : pd.DataFrame
    coluna : str

    Returns
    -------
    pd.DataFrame — linhas cujo valor está fora dos limites de Tukey
    """
    if coluna not in df.columns:
        raise ValueError(f"Coluna '{coluna}' não encontrada.")
    quartis = calcular_quartis(df, coluna)
    q1, q3, iqr = quartis["Q1"], quartis["Q3"], quartis["IQR"]
    # Calcula o limite inferior e superior da regra de Tukey
    limite_inf = q1 - 1.5 * iqr
    limite_sup = q3 + 1.5 * iqr
    # Filtra as linhas que estão fora dos limites
    mascara = (df[coluna] < limite_inf) | (df[coluna] > limite_sup)
    return df[mascara].copy()   # .copy() evita SettingWithCopyWarning


# ── Função 20 ───────────────────────────────────────────────────────────────
def calcular_skewness(df: pd.DataFrame, coluna: str) -> float:
    """
    Calcula o coeficiente de assimetria (skewness) de Fisher.
    Indica se a distribuição tem cauda mais longa à direita (>0)
    ou à esquerda (<0). Valor 0 indica simetria perfeita.

    Parameters
    ----------
    df     : pd.DataFrame
    coluna : str

    Returns
    -------
    float — coeficiente de assimetria
    """
    if coluna not in df.columns:
        raise ValueError(f"Coluna '{coluna}' não encontrada.")
    if not pd.api.types.is_numeric_dtype(df[coluna]):
        raise TypeError(f"Coluna '{coluna}' não é numérica.")
    # skew() do Pandas usa a fórmula ajustada de Fisher (bias=False internamente)
    return float(df[coluna].skew())


# ── Função 21 ───────────────────────────────────────────────────────────────
def calcular_kurtosis(df: pd.DataFrame, coluna: str) -> float:
    """
    Calcula a curtose (excesso de curtose de Fisher) de uma coluna.
    Acima de 0 (leptocúrtica): caudas pesadas, mais outliers.
    Abaixo de 0 (platicúrtica): caudas leves.
    0 = distribuição normal (mesocúrtica).

    Parameters
    ----------
    df     : pd.DataFrame
    coluna : str

    Returns
    -------
    float — excesso de curtose
    """
    if coluna not in df.columns:
        raise ValueError(f"Coluna '{coluna}' não encontrada.")
    if not pd.api.types.is_numeric_dtype(df[coluna]):
        raise TypeError(f"Coluna '{coluna}' não é numérica.")
    # kurt() retorna excesso de curtose (subtrai 3 da fórmula de Pearson)
    return float(df[coluna].kurt())
\end{lstlisting}

% ---------------------------------------------------------------------------
\subsection{Módulo \texttt{limpeza.py} — Nulos e Imputação}
% ---------------------------------------------------------------------------

\begin{lstlisting}[language=Python]
# src/aedtools/limpeza.py
"""
Diagnóstico e tratamento de dados ausentes (missing values).
"""

import pandas as pd


# ── Função 7 ────────────────────────────────────────────────────────────────
def contar_nulos(df: pd.DataFrame) -> pd.DataFrame:
    """
    Retorna DataFrame com contagem absoluta de nulos por coluna,
    valor diferente de proporcao_nulos que retorna percentuais.

    Parameters
    ----------
    df : pd.DataFrame

    Returns
    -------
    pd.DataFrame — colunas: ["coluna", "nulos", "total", "pct"]
    """
    # isnull().sum() conta os NaN de cada coluna
    nulos = df.isnull().sum()
    total = len(df)
    resultado = pd.DataFrame({
        "coluna": nulos.index,
        "nulos":  nulos.values,
        "total":  total,
        # Arredonda para 2 casas decimais para legibilidade
        "pct":    (nulos.values / total * 100).round(2),
    })
    # Ordena do mais nulo para o menos nulo
    return resultado.sort_values("nulos", ascending=False).reset_index(drop=True)


# ── Função 8 ────────────────────────────────────────────────────────────────
def preencher_nulos_media(df: pd.DataFrame, coluna: str) -> pd.DataFrame:
    """
    Imputa NaN de uma coluna numérica com a média da própria coluna.
    Use apenas quando a distribuição é simétrica e sem outliers graves,
    pois a média é sensível a extremos.

    Parameters
    ----------
    df     : pd.DataFrame
    coluna : str — coluna numérica a ser imputada

    Returns
    -------
    pd.DataFrame — cópia do DataFrame com NaN substituídos
    """
    if coluna not in df.columns:
        raise ValueError(f"Coluna '{coluna}' não encontrada.")
    if not pd.api.types.is_numeric_dtype(df[coluna]):
        raise TypeError(f"Coluna '{coluna}' não é numérica.")
    df_out = df.copy()           # nunca modifica o DataFrame original
    media = df_out[coluna].mean()
    # fillna substitui todos os NaN pelo valor informado
    df_out[coluna] = df_out[coluna].fillna(media)
    return df_out


# ── Função 9 ────────────────────────────────────────────────────────────────
def preencher_nulos_mediana(df: pd.DataFrame, coluna: str) -> pd.DataFrame:
    """
    Imputa NaN com a mediana da coluna — mais robusta a outliers
    do que a imputação pela média. Preferível quando há assimetria.

    Parameters
    ----------
    df     : pd.DataFrame
    coluna : str

    Returns
    -------
    pd.DataFrame — cópia com NaN substituídos pela mediana
    """
    if coluna not in df.columns:
        raise ValueError(f"Coluna '{coluna}' não encontrada.")
    if not pd.api.types.is_numeric_dtype(df[coluna]):
        raise TypeError(f"Coluna '{coluna}' não é numérica.")
    df_out = df.copy()
    mediana = df_out[coluna].median()
    df_out[coluna] = df_out[coluna].fillna(mediana)
    return df_out


# ── Função 10 ───────────────────────────────────────────────────────────────
def remover_linhas_nulas(df: pd.DataFrame) -> pd.DataFrame:
    """
    Remove todas as linhas que tenham ao menos um NaN em qualquer coluna.
    Útil para MCAR com baixo percentual de nulos (<5%).
    Não use se houver muitos nulos — perda de dados relevante.

    Parameters
    ----------
    df : pd.DataFrame

    Returns
    -------
    pd.DataFrame — cópia sem linhas com NaN
    """
    # dropna() remove linha se qualquer coluna tiver NaN (how='any' padrão)
    return df.dropna().reset_index(drop=True)


# ── Função 26 ───────────────────────────────────────────────────────────────
def proporcao_nulos(df: pd.DataFrame) -> pd.Series:
    """
    Retorna a proporção (0 a 1) de valores nulos em cada coluna.
    Versão compacta de contar_nulos — útil para filtragem rápida.

    Parameters
    ----------
    df : pd.DataFrame

    Returns
    -------
    pd.Series — índice = nome da coluna, valor = proporção de nulos
    """
    # isnull().mean() calcula a proporção pois True=1 e False=0
    return df.isnull().mean().sort_values(ascending=False)
\end{lstlisting}

% ---------------------------------------------------------------------------
\subsection{Módulo \texttt{transformacao.py} — Normalização, Encoding e Séries Temporais}
% ---------------------------------------------------------------------------

\begin{lstlisting}[language=Python]
# src/aedtools/transformacao.py
"""
Transformações de escala (normalização, z-score), codificação de
variáveis categóricas e cálculos em séries temporais.
"""

import pandas as pd
import numpy as np


# ── Função 14 ───────────────────────────────────────────────────────────────
def normalizar_min_max(df: pd.DataFrame, coluna: str) -> pd.DataFrame:
    """
    Normaliza uma coluna para o intervalo [0, 1] por Min-Max Scaling.
    Fórmula:  x_norm = (x - x_min) / (x_max - x_min)
    Preserva o formato da distribuição original.

    Parameters
    ----------
    df     : pd.DataFrame
    coluna : str — coluna a normalizar

    Returns
    -------
    pd.DataFrame — cópia com coluna '{coluna}_norm' adicionada
    """
    if coluna not in df.columns:
        raise ValueError(f"Coluna '{coluna}' não encontrada.")
    if not pd.api.types.is_numeric_dtype(df[coluna]):
        raise TypeError(f"Coluna '{coluna}' não é numérica.")
    df_out = df.copy()
    x_min = df_out[coluna].min()
    x_max = df_out[coluna].max()
    if x_max == x_min:
        # Todos os valores iguais: normalização não é definida; retorna zeros
        df_out[f"{coluna}_norm"] = 0.0
    else:
        df_out[f"{coluna}_norm"] = (df_out[coluna] - x_min) / (x_max - x_min)
    return df_out


# ── Função 15 ───────────────────────────────────────────────────────────────
def padronizar_zscore(df: pd.DataFrame, coluna: str) -> pd.DataFrame:
    """
    Padroniza uma coluna para média 0 e desvio padrão 1 (z-score).
    Fórmula:  z = (x - média) / DP
    Essencial quando algoritmos de ML são sensíveis à escala.

    Parameters
    ----------
    df     : pd.DataFrame
    coluna : str

    Returns
    -------
    pd.DataFrame — cópia com coluna '{coluna}_z' adicionada
    """
    if coluna not in df.columns:
        raise ValueError(f"Coluna '{coluna}' não encontrada.")
    if not pd.api.types.is_numeric_dtype(df[coluna]):
        raise TypeError(f"Coluna '{coluna}' não é numérica.")
    df_out = df.copy()
    media = df_out[coluna].mean()
    desvio = df_out[coluna].std(ddof=1)   # desvio amostral
    if desvio == 0:
        df_out[f"{coluna}_z"] = 0.0       # todas as obs. iguais
    else:
        df_out[f"{coluna}_z"] = (df_out[coluna] - media) / desvio
    return df_out


# ── Função 22 ───────────────────────────────────────────────────────────────
def contar_valores_unicos(df: pd.DataFrame, coluna: str) -> int:
    """
    Retorna a cardinalidade de uma coluna: quantos valores distintos existem.
    Útil para decidir se uma coluna numérica deve ser tratada como categórica.

    Parameters
    ----------
    df     : pd.DataFrame
    coluna : str

    Returns
    -------
    int — número de valores únicos (excluindo NaN)
    """
    if coluna not in df.columns:
        raise ValueError(f"Coluna '{coluna}' não encontrada.")
    # nunique() conta valores únicos ignorando NaN por padrão
    return int(df[coluna].nunique())


# ── Função 23 ───────────────────────────────────────────────────────────────
def calcular_frequencia(df: pd.DataFrame, coluna: str) -> pd.DataFrame:
    """
    Gera tabela de frequência absoluta, relativa e percentual de uma coluna.

    Parameters
    ----------
    df     : pd.DataFrame
    coluna : str — qualquer tipo (numérico ou categórico)

    Returns
    -------
    pd.DataFrame — colunas: [coluna, "freq_abs", "freq_rel", "pct"]
                   ordenado da categoria mais frequente para a menos frequente
    """
    if coluna not in df.columns:
        raise ValueError(f"Coluna '{coluna}' não encontrada.")
    # value_counts() retorna frequência absoluta ordenada
    freq_abs = df[coluna].value_counts()
    total = freq_abs.sum()
    freq_rel = freq_abs / total            # proporção (0 a 1)
    pct      = freq_rel * 100             # porcentagem (0 a 100)
    resultado = pd.DataFrame({
        coluna:     freq_abs.index,
        "freq_abs": freq_abs.values,
        "freq_rel": freq_rel.values.round(4),
        "pct":      pct.values.round(2),
    })
    return resultado.reset_index(drop=True)


# ── Função 24 ───────────────────────────────────────────────────────────────
def converter_para_categorico(df: pd.DataFrame, coluna: str) -> pd.DataFrame:
    """
    Converte uma coluna para o tipo pd.Categorical do Pandas.
    Reduz uso de memória e habilita ordenamento personalizado.

    Parameters
    ----------
    df     : pd.DataFrame
    coluna : str

    Returns
    -------
    pd.DataFrame — cópia com coluna convertida para Categorical
    """
    if coluna not in df.columns:
        raise ValueError(f"Coluna '{coluna}' não encontrada.")
    df_out = df.copy()
    # pd.Categorical infere as categorias automaticamente
    df_out[coluna] = pd.Categorical(df_out[coluna])
    return df_out


# ── Função 25 ───────────────────────────────────────────────────────────────
def one_hot_encoding(df: pd.DataFrame, coluna: str) -> pd.DataFrame:
    """
    Aplica codificação one-hot (variáveis dummy) em uma coluna categórica.
    Cada categoria vira nova coluna binária (0 ou 1).
    drop_first=True evita multicolinearidade em modelos de regressão.

    Parameters
    ----------
    df     : pd.DataFrame
    coluna : str — coluna categórica a codificar

    Returns
    -------
    pd.DataFrame — DataFrame com a coluna original removida e novas
                   colunas binárias adicionadas
    """
    if coluna not in df.columns:
        raise ValueError(f"Coluna '{coluna}' não encontrada.")
    df_out = df.copy()
    # get_dummies cria colunas booleanas para cada categoria
    dummies = pd.get_dummies(df_out[coluna], prefix=coluna, drop_first=True)
    # Remove coluna original e concatena as dummies
    df_out = df_out.drop(columns=[coluna])
    df_out = pd.concat([df_out, dummies], axis=1)
    return df_out


# ── Função 27 ───────────────────────────────────────────────────────────────
def soma_acumulada(df: pd.DataFrame, coluna: str) -> pd.DataFrame:
    """
    Calcula a soma cumulativa de uma coluna linha por linha.
    Útil para saldo ao longo do tempo ou estoque acumulado.

    Parameters
    ----------
    df     : pd.DataFrame
    coluna : str — coluna numérica

    Returns
    -------
    pd.DataFrame — cópia com coluna '{coluna}_acum' adicionada
    """
    if coluna not in df.columns:
        raise ValueError(f"Coluna '{coluna}' não encontrada.")
    if not pd.api.types.is_numeric_dtype(df[coluna]):
        raise TypeError(f"Coluna '{coluna}' não é numérica.")
    df_out = df.copy()
    # cumsum() soma todos os valores anteriores mais o atual
    df_out[f"{coluna}_acum"] = df_out[coluna].cumsum()
    return df_out


# ── Função 28 ───────────────────────────────────────────────────────────────
def media_movel(df: pd.DataFrame, coluna: str, janela: int = 7) -> pd.DataFrame:
    """
    Calcula a média móvel simples com janela configurável.
    Suaviza flutuações de curto prazo em séries temporais.

    Parameters
    ----------
    df     : pd.DataFrame
    coluna : str — série temporal numérica
    janela : int — tamanho da janela (default: 7 períodos)

    Returns
    -------
    pd.DataFrame — cópia com coluna '{coluna}_mm{janela}' adicionada
                   (NaN nas primeiras janela-1 linhas é esperado)
    """
    if coluna not in df.columns:
        raise ValueError(f"Coluna '{coluna}' não encontrada.")
    if not pd.api.types.is_numeric_dtype(df[coluna]):
        raise TypeError(f"Coluna '{coluna}' não é numérica.")
    if janela < 1:
        raise ValueError("Janela deve ser >= 1.")
    df_out = df.copy()
    # rolling(janela).mean() calcula a média dos últimos 'janela' valores
    df_out[f"{coluna}_mm{janela}"] = df_out[coluna].rolling(window=janela).mean()
    return df_out


# ── Função 29 ───────────────────────────────────────────────────────────────
def variacao_percentual(df: pd.DataFrame, coluna: str) -> pd.DataFrame:
    """
    Calcula a variação percentual entre linhas consecutivas.
    Fórmula: (valor_t - valor_{t-1}) / valor_{t-1} * 100
    Útil para séries de preços, receitas ou qualquer métrica temporal.

    Parameters
    ----------
    df     : pd.DataFrame
    coluna : str

    Returns
    -------
    pd.DataFrame — cópia com coluna '{coluna}_var_pct' adicionada
                   (primeira linha será NaN pois não há anterior)
    """
    if coluna not in df.columns:
        raise ValueError(f"Coluna '{coluna}' não encontrada.")
    if not pd.api.types.is_numeric_dtype(df[coluna]):
        raise TypeError(f"Coluna '{coluna}' não é numérica.")
    df_out = df.copy()
    # pct_change() calcula a variação % em relação ao valor anterior
    df_out[f"{coluna}_var_pct"] = df_out[coluna].pct_change() * 100
    return df_out


# ── Função 30 ───────────────────────────────────────────────────────────────
def calcular_ranking(
    df: pd.DataFrame, coluna: str, crescente: bool = False
) -> pd.DataFrame:
    """
    Adiciona colunas de ranking crescente e decrescente a partir de uma coluna.
    Empates recebem o mesmo rank (method='min' — comportamento esportivo).

    Parameters
    ----------
    df        : pd.DataFrame
    coluna    : str — coluna numérica para rankear
    crescente : bool — se True o rank 1 é o menor valor;
                       se False (padrão) o rank 1 é o maior valor

    Returns
    -------
    pd.DataFrame — cópia com coluna '{coluna}_rank' adicionada
    """
    if coluna not in df.columns:
        raise ValueError(f"Coluna '{coluna}' não encontrada.")
    if not pd.api.types.is_numeric_dtype(df[coluna]):
        raise TypeError(f"Coluna '{coluna}' não é numérica.")
    df_out = df.copy()
    # ascending=True -> menor valor recebe rank 1
    # ascending=False -> maior valor recebe rank 1 (ranking de desempenho)
    df_out[f"{coluna}_rank"] = df_out[coluna].rank(
        ascending=crescente, method="min"
    ).astype(int)
    return df_out


# ── Função 31 ───────────────────────────────────────────────────────────────
def identificar_tipos_colunas(df: pd.DataFrame) -> dict:
    """
    Classifica as colunas do DataFrame em numéricas e categóricas.
    Útil como primeiro passo de qualquer pipeline de AED.

    Parameters
    ----------
    df : pd.DataFrame

    Returns
    -------
    dict — chaves: "numericas" (list[str]), "categoricas" (list[str])
    """
    # select_dtypes filtra colunas por família de tipo
    numericas   = df.select_dtypes(include=["number"]).columns.tolist()
    categoricas = df.select_dtypes(
        include=["object", "category", "bool"]
    ).columns.tolist()
    return {"numericas": numericas, "categoricas": categoricas}
\end{lstlisting}

% ---------------------------------------------------------------------------
\subsection{Módulo \texttt{correlacao.py} — Correlação, Métricas e Intervalos}
% ---------------------------------------------------------------------------

\begin{lstlisting}[language=Python]
# src/aedtools/correlacao.py
"""
Correlação de Pearson, matriz de correlação e métricas de avaliação
de modelos preditivos (R², MAE, MSE, RMSE) e intervalo de confiança.
"""

import pandas as pd
import numpy as np
from scipy import stats


# ── Função 16 ───────────────────────────────────────────────────────────────
def calcular_correlacao(
    df: pd.DataFrame, col1: str, col2: str
) -> float:
    """
    Calcula o coeficiente de correlação de Pearson entre duas colunas.
    Mede apenas relação LINEAR. Use Spearman para variáveis ordinais.
    Escala: -1 (negativa perfeita) a +1 (positiva perfeita). 0 = sem correlação linear.

    Parameters
    ----------
    df   : pd.DataFrame
    col1 : str — primeira coluna numérica
    col2 : str — segunda coluna numérica

    Returns
    -------
    float — r de Pearson
    """
    for col in [col1, col2]:
        if col not in df.columns:
            raise ValueError(f"Coluna '{col}' não encontrada.")
        if not pd.api.types.is_numeric_dtype(df[col]):
            raise TypeError(f"Coluna '{col}' não é numérica.")
    # corr() ignora automaticamente as linhas com NaN em qualquer das colunas
    return float(df[col1].corr(df[col2]))


# ── Função 32 ───────────────────────────────────────────────────────────────
def gerar_matriz_correlacao(df: pd.DataFrame) -> pd.DataFrame:
    """
    Gera a matriz de correlação de Pearson para todas as colunas numéricas.
    Cada célula [i,j] é o r de Pearson entre a coluna i e a coluna j.
    A diagonal principal é sempre 1.0 (correlação de uma variável consigo mesma).

    Parameters
    ----------
    df : pd.DataFrame

    Returns
    -------
    pd.DataFrame — matriz quadrada de correlações (linhas = colunas = atributos)
    """
    # Seleciona apenas colunas numéricas antes de calcular
    df_num = df.select_dtypes(include=["number"])
    if df_num.empty:
        raise ValueError("Nenhuma coluna numérica encontrada no DataFrame.")
    # corr() aplica Pearson a cada par de colunas
    return df_num.corr(method="pearson")


# ── Função 34 ───────────────────────────────────────────────────────────────
def calcular_r2(y_real: pd.Series, y_pred: pd.Series) -> float:
    """
    Coeficiente de Determinação R² — mede quanto da variância
    em y_real é explicada pelo modelo (y_pred).
    R²=1: ajuste perfeito. R²=0: modelo não melhor que a média.
    R²<0: modelo pior que a média (possível overfitting ou dados ruins).

    Parameters
    ----------
    y_real : pd.Series — valores reais observados
    y_pred : pd.Series — valores preditos pelo modelo

    Returns
    -------
    float — R² (pode ser negativo)
    """
    y_real = np.array(y_real)
    y_pred = np.array(y_pred)
    # SS_res: soma dos quadrados dos resíduos (erro do modelo)
    ss_res = np.sum((y_real - y_pred) ** 2)
    # SS_tot: soma dos quadrados totais (variância total de y_real)
    ss_tot = np.sum((y_real - np.mean(y_real)) ** 2)
    if ss_tot == 0:
        return 1.0  # todos os valores reais são iguais: R² = 1 por convenção
    return float(1 - ss_res / ss_tot)


# ── Função 35 ───────────────────────────────────────────────────────────────
def calcular_mae(y_real: pd.Series, y_pred: pd.Series) -> float:
    """
    Erro Absoluto Médio (MAE — Mean Absolute Error).
    Mede o erro médio sem penalizar erros grandes — mais intuitivo.
    Unidade: a mesma de y_real.

    Parameters
    ----------
    y_real : pd.Series
    y_pred : pd.Series

    Returns
    -------
    float — MAE
    """
    y_real = np.array(y_real)
    y_pred = np.array(y_pred)
    # Diferença absoluta linha a linha, depois média
    return float(np.mean(np.abs(y_real - y_pred)))


# ── Função 36 ───────────────────────────────────────────────────────────────
def calcular_mse(y_real: pd.Series, y_pred: pd.Series) -> float:
    """
    Erro Quadrático Médio (MSE — Mean Squared Error).
    Penaliza erros grandes mais do que pequenos (eleva ao quadrado).
    Unidade: quadrado da unidade de y_real — menos intuitivo que MAE.

    Parameters
    ----------
    y_real : pd.Series
    y_pred : pd.Series

    Returns
    -------
    float — MSE
    """
    y_real = np.array(y_real)
    y_pred = np.array(y_pred)
    # Eleva o resíduo ao quadrado e calcula a média
    return float(np.mean((y_real - y_pred) ** 2))


# ── Função 37 ───────────────────────────────────────────────────────────────
def calcular_rmse(y_real: pd.Series, y_pred: pd.Series) -> float:
    """
    Raiz do Erro Quadrático Médio (RMSE — Root Mean Squared Error).
    Combina a penalidade dos erros grandes do MSE com a interpretabilidade
    do MAE: está na mesma unidade de y_real.

    Parameters
    ----------
    y_real : pd.Series
    y_pred : pd.Series

    Returns
    -------
    float — RMSE
    """
    # Simplesmente a raiz quadrada do MSE
    return float(np.sqrt(calcular_mse(y_real, y_pred)))


# ── Função 38 ───────────────────────────────────────────────────────────────
def intervalo_confianca_media(
    df: pd.DataFrame, coluna: str, alpha: float = 0.05
) -> dict:
    """
    Calcula o intervalo de confiança para a média usando a distribuição t.
    Para amostras grandes (n>30) converge para o IC normal. Para n<30 o
    uso do t-Student é obrigatório.

    Fórmula:  IC = x̄  ±  t_{α/2, n-1}  ×  (s / √n)

    Parameters
    ----------
    df     : pd.DataFrame
    coluna : str — coluna numérica
    alpha  : float — nível de significância (default 0.05 → IC de 95%)

    Returns
    -------
    dict — chaves: "media", "limite_inf", "limite_sup", "confianca_pct"
    """
    if coluna not in df.columns:
        raise ValueError(f"Coluna '{coluna}' não encontrada.")
    if not pd.api.types.is_numeric_dtype(df[coluna]):
        raise TypeError(f"Coluna '{coluna}' não é numérica.")
    serie = df[coluna].dropna()
    n     = len(serie)
    if n < 2:
        raise ValueError("São necessárias ao menos 2 observações.")
    media  = serie.mean()
    erro_p = serie.sem()   # erro padrão da média = s / sqrt(n)
    # t.ppf: quantil da distribuição t com graus de liberdade n-1
    t_crit = stats.t.ppf(1 - alpha / 2, df=n - 1)
    margem = t_crit * erro_p
    return {
        "media":         round(float(media), 4),
        "limite_inf":    round(float(media - margem), 4),
        "limite_sup":    round(float(media + margem), 4),
        "confianca_pct": round((1 - alpha) * 100, 1),
    }
\end{lstlisting}

% ---------------------------------------------------------------------------
\subsection{Módulo \texttt{visualizacao.py} — Gráficos com Matplotlib e Seaborn}
% ---------------------------------------------------------------------------

\begin{lstlisting}[language=Python]
# src/aedtools/visualizacao.py
"""
Funções de visualização: histograma, boxplot, scatter, heatmap,
barras e pizza. Todas retornam o objeto fig para personalização
posterior e evitam plt.show() automático (deixa o usuário decidir).
"""

import pandas as pd
import numpy as np
import matplotlib.pyplot as plt
import seaborn as sns


# ── Função 17 ───────────────────────────────────────────────────────────────
def plotar_histograma(
    df: pd.DataFrame, coluna: str, bins: int = 20,
    titulo: str | None = None
) -> plt.Figure:
    """
    Plota histograma com curva KDE de uma coluna numérica.
    Inclui linhas verticais de média e mediana para comparação visual.

    Parameters
    ----------
    df      : pd.DataFrame
    coluna  : str
    bins    : int — número de intervalos (default 20)
    titulo  : str | None — título do gráfico (auto-gerado se None)

    Returns
    -------
    matplotlib.figure.Figure
    """
    if coluna not in df.columns:
        raise ValueError(f"Coluna '{coluna}' não encontrada.")
    if not pd.api.types.is_numeric_dtype(df[coluna]):
        raise TypeError(f"Coluna '{coluna}' não é numérica.")
    serie = df[coluna].dropna()
    fig, ax = plt.subplots(figsize=(8, 4))
    # histplot com kde=True sobrepõe curva de densidade suavizada
    sns.histplot(serie, bins=bins, kde=True, ax=ax, color="steelblue", alpha=0.7)
    # Linha vertical da média (tracejada vermelha)
    ax.axvline(serie.mean(), color="red", linestyle="--", linewidth=1.5,
               label=f"Média: {serie.mean():.2f}")
    # Linha vertical da mediana (tracejada verde)
    ax.axvline(serie.median(), color="green", linestyle="--", linewidth=1.5,
               label=f"Mediana: {serie.median():.2f}")
    ax.set_xlabel(coluna)
    ax.set_ylabel("Contagem")
    ax.set_title(titulo or f"Distribuição de {coluna}")
    ax.legend()
    # tight_layout evita corte de labels
    fig.tight_layout()
    return fig


# ── Função 18 ───────────────────────────────────────────────────────────────
def plotar_boxplot(
    df: pd.DataFrame, coluna: str,
    grupo: str | None = None,
    titulo: str | None = None
) -> plt.Figure:
    """
    Plota boxplot de uma coluna numérica, opcionalmente agrupado
    por uma coluna categórica. Outliers são marcados como pontos.

    Parameters
    ----------
    df      : pd.DataFrame
    coluna  : str — variável numérica
    grupo   : str | None — coluna categórica para split (opcional)
    titulo  : str | None

    Returns
    -------
    matplotlib.figure.Figure
    """
    if coluna not in df.columns:
        raise ValueError(f"Coluna '{coluna}' não encontrada.")
    fig, ax = plt.subplots(figsize=(8, 4))
    # Se grupo informado, cria boxplot por categoria; senão, único boxplot
    if grupo and grupo in df.columns:
        sns.boxplot(data=df, x=grupo, y=coluna, ax=ax,
                    palette="pastel", flierprops={"marker": "o"})
        ax.set_title(titulo or f"Boxplot de {coluna} por {grupo}")
    else:
        sns.boxplot(data=df, y=coluna, ax=ax,
                    color="steelblue", flierprops={"marker": "o"})
        ax.set_title(titulo or f"Boxplot de {coluna}")
    fig.tight_layout()
    return fig


# ── Função 19 ───────────────────────────────────────────────────────────────
def plotar_scatterplot(
    df: pd.DataFrame, col1: str, col2: str,
    hue: str | None = None,
    titulo: str | None = None
) -> plt.Figure:
    """
    Scatter plot entre dois atributos numéricos com linha de tendência (regressão).
    Parâmetro hue permite colorir pontos por uma terceira variável categórica.

    Parameters
    ----------
    df     : pd.DataFrame
    col1   : str — variável no eixo X
    col2   : str — variável no eixo Y
    hue    : str | None — coluna categórica para cor dos pontos (opcional)
    titulo : str | None

    Returns
    -------
    matplotlib.figure.Figure
    """
    for col in [col1, col2]:
        if col not in df.columns:
            raise ValueError(f"Coluna '{col}' não encontrada.")
    fig, ax = plt.subplots(figsize=(7, 5))
    # scatterplot com hue opcional
    sns.scatterplot(data=df, x=col1, y=col2, hue=hue,
                    ax=ax, alpha=0.7, edgecolor="white")
    # regplot adiciona linha de tendência linear sem sobrescrever os pontos
    sns.regplot(data=df, x=col1, y=col2, scatter=False,
                ax=ax, color="red", line_kws={"linewidth": 1.5})
    # Calcular e exibir r de Pearson no título
    r = df[col1].corr(df[col2])
    ax.set_title(titulo or f"{col1} vs {col2}  (r = {r:.3f})")
    ax.set_xlabel(col1)
    ax.set_ylabel(col2)
    fig.tight_layout()
    return fig


# ── Função 33 ───────────────────────────────────────────────────────────────
def plotar_heatmap_correlacao(df: pd.DataFrame, titulo: str | None = None) -> plt.Figure:
    """
    Plota heatmap colorido da matriz de correlação para todas as colunas numéricas.
    Usa diverging colormap (azul=negativo, branco=0, vermelho=positivo).
    Anota cada célula com o valor de r arredondado a 2 casas.

    Parameters
    ----------
    df     : pd.DataFrame
    titulo : str | None

    Returns
    -------
    matplotlib.figure.Figure
    """
    # Importa a função do próprio módulo correlacao para consistência
    from .correlacao import gerar_matriz_correlacao
    matriz = gerar_matriz_correlacao(df)
    n = len(matriz)
    # Ajusta o tamanho da figura de acordo com o número de colunas
    tamanho = max(6, n * 0.8)
    fig, ax = plt.subplots(figsize=(tamanho, tamanho))
    # annot=True: escreve o valor em cada célula; fmt=".2f": 2 casas decimais
    sns.heatmap(
        matriz, annot=True, fmt=".2f", ax=ax,
        cmap="coolwarm", center=0,
        vmin=-1, vmax=1,
        linewidths=0.5, linecolor="white",
        square=True,
    )
    ax.set_title(titulo or "Matriz de Correlação de Pearson")
    fig.tight_layout()
    return fig


# ── Função 39 ───────────────────────────────────────────────────────────────
def plotar_barras_categorias(
    df: pd.DataFrame, coluna: str,
    orientacao: str = "v",
    titulo: str | None = None
) -> plt.Figure:
    """
    Gráfico de barras com contagem por categoria, ordenado da mais
    frequente para a menos frequente. Inclui rótulo numérico em cada barra.

    Parameters
    ----------
    df          : pd.DataFrame
    coluna      : str — coluna categórica
    orientacao  : str — 'v' (vertical) ou 'h' (horizontal, melhor para labels longos)
    titulo      : str | None

    Returns
    -------
    matplotlib.figure.Figure
    """
    if coluna not in df.columns:
        raise ValueError(f"Coluna '{coluna}' não encontrada.")
    # value_counts() já retorna ordenado do mais frequente
    contagem = df[coluna].value_counts()
    fig, ax = plt.subplots(figsize=(8, 4))
    if orientacao == "h":
        # Barras horizontais — melhor para categorias com nomes longos
        barras = ax.barh(contagem.index, contagem.values, color="steelblue")
        ax.set_xlabel("Contagem")
        # Rótulo numérico ao lado de cada barra
        for b in barras:
            ax.text(b.get_width() + 0.1, b.get_y() + b.get_height() / 2,
                    str(int(b.get_width())), va="center", fontsize=9)
    else:
        # Barras verticais — padrão
        barras = ax.bar(contagem.index, contagem.values, color="steelblue")
        ax.set_ylabel("Contagem")
        # Rótulo numérico acima de cada barra
        for b in barras:
            ax.text(b.get_x() + b.get_width() / 2, b.get_height() + 0.1,
                    str(int(b.get_height())), ha="center", fontsize=9)
    ax.set_xlabel(coluna)
    ax.set_title(titulo or f"Frequência por {coluna}")
    fig.tight_layout()
    return fig


# ── Função 40 ───────────────────────────────────────────────────────────────
def plotar_pizza_categorias(
    df: pd.DataFrame, coluna: str,
    titulo: str | None = None,
    max_fatias: int = 6
) -> plt.Figure:
    """
    Gráfico de pizza com proporção por categoria.
    Limita a max_fatias categorias (agrupa o resto em 'Outros')
    para evitar visualmente mais de 6 fatias — limite cognitivo humano.

    Parameters
    ----------
    df         : pd.DataFrame
    coluna     : str — coluna categórica
    titulo     : str | None
    max_fatias : int — máximo de categorias exibidas individualmente (default 6)

    Returns
    -------
    matplotlib.figure.Figure
    """
    if coluna not in df.columns:
        raise ValueError(f"Coluna '{coluna}' não encontrada.")
    contagem = df[coluna].value_counts()
    # Se há mais que max_fatias categorias, agrupa o resto em "Outros"
    if len(contagem) > max_fatias:
        top   = contagem.iloc[:max_fatias]
        outros = contagem.iloc[max_fatias:].sum()
        contagem = pd.concat([top, pd.Series([outros], index=["Outros"])])
    fig, ax = plt.subplots(figsize=(6, 6))
    # autopct: exibe a percentagem dentro de cada fatia
    # startangle: começa do topo (90°) como calendário
    ax.pie(
        contagem.values,
        labels=contagem.index,
        autopct="%1.1f%%",
        startangle=90,
        pctdistance=0.82,
        colors=sns.color_palette("pastel", n_colors=len(contagem)),
    )
    ax.set_title(titulo or f"Proporção por {coluna}")
    fig.tight_layout()
    return fig
\end{lstlisting}

\subsection*{Teste unitário — exemplo com \texttt{pytest}}

\begin{lstlisting}[language=Python]
# tests/test_estatistica.py
"""
Testes unitários para o módulo estatistica.py.
Execute com: pytest tests/ -v
"""

import pytest
import pandas as pd
import numpy as np
from aedtools import (
    calcular_media, calcular_mediana, calcular_moda,
    calcular_amplitude, detectar_outliers_iqr,
    calcular_quartis,
)


# Fixture: DataFrame reutilizado em todos os testes deste módulo
@pytest.fixture
def df_simples():
    """Dataset mínimo com valores conhecidos para testes determinísticos."""
    return pd.DataFrame({
        "valores": [1.0, 2.0, 3.0, 4.0, 100.0],  # 100 é outlier intencional
        "categoria": ["A", "B", "A", "C", "A"],
    })


def test_calcular_media(df_simples):
    """Média de [1,2,3,4,100] deve ser 22.0"""
    assert calcular_media(df_simples, "valores") == pytest.approx(22.0)


def test_calcular_mediana(df_simples):
    """Mediana de [1,2,3,4,100] deve ser 3.0"""
    assert calcular_mediana(df_simples, "valores") == pytest.approx(3.0)


def test_calcular_moda_categorico(df_simples):
    """Moda da categoria deve ser 'A' (aparece 3x)"""
    assert calcular_moda(df_simples, "categoria") == ["A"]


def test_calcular_amplitude(df_simples):
    """Amplitude de [1,2,3,4,100] deve ser 99"""
    assert calcular_amplitude(df_simples, "valores") == pytest.approx(99.0)


def test_detectar_outliers_iqr(df_simples):
    """100 deve ser detectado como outlier pela regra de Tukey"""
    outliers = detectar_outliers_iqr(df_simples, "valores")
    assert 100.0 in outliers["valores"].values
    assert len(outliers) == 1


def test_coluna_inexistente_levanta_valueerror(df_simples):
    """Passar coluna que não existe deve levantar ValueError"""
    with pytest.raises(ValueError, match="não encontrada"):
        calcular_media(df_simples, "nao_existe")


def test_funcao_nao_modifica_original(df_simples):
    """Funções não devem modificar o DataFrame original (imutabilidade)"""
    df_original = df_simples.copy()
    calcular_amplitude(df_simples, "valores")
    pd.testing.assert_frame_equal(df_simples, df_original)
\end{lstlisting}

\begin{TipBox}
\textbf{Dica de workflow para a turma:}
\begin{enumerate}[leftmargin=*]
  \item Clone o repositório: \texttt{git clone https://github.com/unifacol/aedtools}
  \item Crie seu branch: \texttt{git checkout -b feat/funcao-XX-seu-nome}
  \item Implemente a sua função no módulo correto.
  \item Adicione pelo menos 3 testes em \texttt{tests/}.
  \item Abra um Pull Request — o professor revisa e aprova o merge.
  \item Após aprovação, a função entra oficialmente na \texttt{v0.X.0} da biblioteca.
\end{enumerate}
Ao final do semestre, a biblioteca tem 40 funções testadas, documentadas e
publicadas no PyPI — um portfólio real para o LinkedIn e o currículo.
\end{TipBox}
